\section{Introduction}\label{introduction}

In the model presented here, the phenomenon of gravity is not introduced
as a fundamental force, but rather emerges from the internal dynamics of
a tensioned medium --- a three-dimensional brane embedded in a
higher-dimensional space. This use of the term ``brane'' should not be confused with Randall--Sundrum--type brane-world scenarios, in which a thin 3+1-dimensional brane with tuned tension sits in a warped higher-dimensional Einstein-gravity bulk; in the present work, the brane is a physical elastic medium embedded in flat four-dimensional space, and gravity is interpreted as emergent from its induced metric and internal tension dynamics. For intuition we often adopt an embedding frame in which the ground-state brane is approximately aligned with the first three embedding coordinates $(X^1,X^2,X^3)$ and transverse displacement is represented primarily by the remaining coordinate $X^4$. This is a coordinate/gauge convenience only: the substrate dynamics depends on the full embedding vector $\mathbf X\in\mathbb R^4$ and does not privilege any component.

A key feature of this model is that transverse (out-of-brane) deformations are
geometrically coupled to in-brane stretch: because the induced metric
$g_{ij}=\partial_i\mathbf X\cdot\partial_j\mathbf X$ depends on all embedding
components, localized transverse excitation generically modifies in-brane
distances and therefore the local elastic energy. The brane relaxes by adjusting
its in-brane stretches, producing spatially extended lateral contraction
patterns. Energy concentrations --- represented by oscillatory or solitonic
structures with significant transverse deformation --- can therefore induce
curvature-like behavior in the effective geometry experienced by long-wavelength
waves. The resulting contraction fields resemble the attractive behavior
attributed to gravity in general relativity, yet they arise here from local
tension dynamics within a medium governed by deterministic laws. Thus, gravity
appears as a \textbf{secondary, emergent effect} of the system's effort to
minimize internal deformation --- a direct consequence of the geometric coupling
between transverse displacement and in-brane stretch.

A central heuristic in this work is that \emph{mass is nothing but
localized energy in motion}, in line with the standard relativistic
equivalence between inertial and gravitational mass. In particular, the
gravitational mass of any closed system is taken to be proportional to
its total energy in the center-of-mass frame, including kinetic and
internal degrees of freedom. This viewpoint is emphasized clearly by
van der Mark and 't~Hooft in their essay \emph{``Light is Heavy''},
where they show that a box filled with radiation weighs more than an
empty one, and that the gravitational mass of the light inside is
precisely $E/c^2$ even though each photon has zero rest mass
\cite{vanderMark2000LightHeavy}. The present brane model simply
implements this idea mechanically: all mass, whether ``matter'' or
``radiation'', is reinterpreted as internal energy stored in waves and
tension of a single medium.

The present work builds on the toroidal electron model of Williamson and van der Mark (1997), generalizing it into a continuous brane framework where confinement, relativity, and gravitation emerge from a unified geometric substrate (detailed in Section~\ref{sec:internal-structure-of-the-electron}).

\subsection{Core Idea and Scope}\label{core-idea}

The central hypothesis explored here is that \textbf{lateral contraction} in a tensioned medium—typically neglected under the small-angle approximation in classical wave theory—may, when extended to higher dimensions, provide a natural mechanism for gravitational effects. Specifically, if a 3D brane is embedded in 4D space, then localized transverse excitation can, via the induced-metric coupling, drive spatially extended contraction patterns in the in-brane directions that resemble gravitational fields.

This paper develops that idea into a coherent framework in which:
\begin{itemize}
\item \textbf{Special relativity} emerges from isotropic wave propagation in an absolute medium;
\item \textbf{Quantum uncertainty} reflects Fourier-localization constraints on wave packets;
\item \textbf{Particles} correspond to topologically stable solitons with internal phase structure;
\item \textbf{Electromagnetism} arises as an emergent $U(1)$ gauge field from geometric (Berry) phases of band-isolated substrate modes;
\item \textbf{Gravitation} emerges from transverse-deformation-induced curvature of the brane's induced metric.
\end{itemize}

While speculative and exploratory, the framework aims to demonstrate that a single elastic substrate can, in principle, unify phenomena typically treated as separate in standard physics.

\paragraph{Microscopic time vs emergent spacetime.}
Throughout this work we distinguish between a microscopic evolution parameter $t$ and the emergent spacetime experienced by observers. At the fundamental level, the brane configuration
\[
\mathbf X : \Omega \times \mathbb{R} \to \mathbb{R}^4, \qquad (x,t) \mapsto \mathbf X(x,t)
\]
evolves with respect to an external time parameter $t$ that is not part of the embedding itself. At the emergent level, observers reconstruct effective spacetime coordinates $x^\mu = (ct, x^i)$ and an approximate Lorentzian metric $g_{\mu\nu}(x)$ from the propagation of waves on the brane. In other words, relativity arises as an effective description of wave kinematics on a pre-existing brane evolving in microscopic time $t$.

\paragraph{Structure of this paper.}
The paper follows a substrate--first perspective.
Section~\ref{conceptual-model} defines the microscopic brane substrate (degrees of freedom, action, equations of motion)
and introduces the minimal geometric-phase (Berry) construction we will later apply to prepared narrowband wave packets.
Section~\ref{emergent-relativity-and-lorentz-invariance} discusses how Lorentz-like kinematics can arise effectively from the
linear wave regime of the substrate.
Section~\ref{sec:internal-structure-of-the-electron} specifies the target internal electron structure (toroidal/solitonic)
that motivates the later experiments.
Section~\ref{experimental-setting} describes the discrete simulation framework and the preparation of narrowband carriers.
In this revision, the more elaborate geometric-phase reconstructions (curvature maps, defect/source localization, non-Abelian transport)
and any full identification with Maxwell theory are deferred; we proceed step by step from stable initialization toward
minimal geometric phase measurements.

\subsection{Related Works}\label{related-works}

Recent work by \textbf{Roth et al.}~\cite{Roth2023QQM2} and \textbf{Close}~\cite{Close2025EquationEverything}
develops closely related elastic-medium approaches in which wave dynamics in a structured substrate
is taken as fundamental and particle-like objects arise as localized excitations. While these
approaches differ in mathematical details and ontology, they overlap with our motivation to treat
relativistic-looking behavior and field-like phenomena as emergent from underlying medium dynamics.

A second, directly relevant lineage is the theory of \emph{geometric phase} and the emergence of
\emph{gauge structure} from adiabatic transport in a band/eigenspace. The Berry phase~\cite{Berry1984}
and its degenerate/non-Abelian generalization~\cite{WilczekZee1984} formalize how a connection and
holonomy arise from parallel transport in a parameter-dependent mode subspace. In our setting, this
mathematical structure is used as part of the proposed bridge from substrate waves to an effective gauge description: after isolating a narrowband substrate mode, we
compute the associated connection/holonomy to assess whether an effective gauge-like description is
supported by the simulated brane dynamics (see also~\cite{Katanaev2009} for geometric interpretation).

A complementary line of work studies electron and photon structure on Compton scales. The toroidal
model of Williamson and van der Mark treats the electron as a bound circulating excitation and
highlights the special role of $\lambda_C$ in setting a characteristic internal scale~\cite{Williamson1997,Williamson2015}.
Zitterbewegung-inspired approaches likewise connect internal motion radii to the reduced Compton
wavelength~\cite{UrdanetaSantos2023}. These strands inform our choice of Compton-calibrated lattice
scales and our soliton-inspired electron target geometry.